\chapter{Symbolic sums and momentum space}
\setwlchapter{momentum}
\section{An index belongs to its sum}
In $\sum_k a_k^\dagger a_k$ the value of $k$ has no meaning outside
the sum. Multiplying two copies requires two independent summation
indices, even if the same letter was used to write both. Restricting
them to be equal would throw away every term coupling two different
levels. The symbolic head \fn{sum} preserves this distinction;
\fn{Sum} with numerical limits instead enumerates a cluster.
\casestudy{dummies}{Nested summations}

\section{Momentum transfer and Fourier convention}
For $N_s$ Bravais cells take
\[
 c_R=N_s^{-1/2}\sum_k e^{ik\cdot R}c_k,\qquad
 c_k=N_s^{-1/2}\sum_R e^{-ik\cdot R}c_R .
\]
One independent sum supplies one factor $N_s$ when a lattice phase
is constant. A density bilinear brings $1/N_s$ but its site sum brings
$N_s$ back; a local quartic vertex retains $1/N_s$ and three independent
momenta. Declare lattice and momentum indices with the \emph{same}
lattice size. Keep \fn{blochphase} symbolic when combining deltas;
convert to exponentials only for an explicit dispersion or geometry.
\casestudy{transfer}{Following the transfer through a density operator}

\section{Periodic clusters and boundary twists}
In one dimension $c_{n+L}=e^{i\theta}c_n$ selects the grid
$k_m=(2\pi m+\theta)/L$. Bilinears involve differences of two
momenta: the twists cancel in a transferred momentum. Thus it is
incorrect to shift \emph{every} momentum in a momentum-conservation
delta by $\theta/L$. \fn{clusterhop} inserts the phase only when a
real-space hop crosses the boundary; \fn{fourierrules} gives the
explicit finite-dimensional change of basis. In a periodic untwisted
cluster \fn{sumCluster} folds indices modulo each lattice dimension.
\casestudy{twist}{Flux sensitivity of the spectrum}
\begin{exercise}[Discrete unitarity]\label{ex:grid}
Show that a common twist does not change the inner product of two
columns of the finite Fourier matrix. Test $L=3$, $\theta=\pi/2$ and
explain why a reversed sign in just the inverse transform fails.
\end{exercise}
